\documentclass[11pt,a4paper]{article}
\usepackage[margin=19mm,headheight=14pt]{geometry}
\usepackage[T1]{fontenc}
\usepackage{lmodern,amsmath,amssymb,booktabs,tabularx,enumitem,xcolor,fancyhdr,microtype}
\usepackage[hidelinks]{hyperref}
\definecolor{Ink}{HTML}{16324F}
\definecolor{Accent}{HTML}{216B78}
\pagestyle{fancy}
\fancyhf{}
\fancyhead[L]{\small\sffamily Scientific Computing 2026}
\fancyhead[R]{\small\sffamily Lectures 1--7}
\fancyfoot[C]{\small\thepage}
\renewcommand{\headrulewidth}{0.3pt}
\setlength{\parindent}{0pt}
\setlength{\parskip}{5pt}
\setlength{\emergencystretch}{2em}
\newcommand{\blocktitle}[2]{{\Large\sffamily\bfseries\color{Ink}Block #1: #2}\par\vspace{5pt}}
\newenvironment{question}[2]{\par\addvspace{10pt}\begin{minipage}{\linewidth}{\sffamily\bfseries\color{Ink}#1.\ #2}\par\vspace{3pt}}{\end{minipage}\par}
\newcommand{\E}{\mathbb E}
\newcommand{\Var}{\operatorname{Var}}
\newcommand{\Cov}{\operatorname{Cov}}
\hypersetup{pdftitle={Colloquium preparation: Lectures 1--7},pdfauthor={Scientific Computing 2026}}
\begin{document}
{\LARGE\sffamily\bfseries\color{Ink}Colloquium preparation}\par
{\sffamily Lectures 1--7 \hfill 20 minutes per student}
\vspace{3pt}

You will receive three questions chosen at random: one from each of the
three blocks below. You will have 20 minutes in total to answer them orally,
write down your calculations, and respond to follow-up questions.

This is the question bank for preparation. You have until the colloquium
to prepare the questions in all three blocks. Be ready to explain your
reasoning and the meaning of the formulas you use. The calculation problems
are intended to be solved by hand; no coding is required. Logarithms and
exponentials may be left unevaluated.

\blocktitle{A}{Definitions and interpretation}
For your selected question, explain the concepts and write the relevant
formulas. No proofs are required.

{\small
\begin{enumerate}[label=\textbf{A\arabic*.},leftmargin=10mm,itemsep=4pt,topsep=3pt,parsep=0pt]
\item What is a random variable? What is its cumulative distribution function?
\item What is conditional probability? What does it mean for two events to be independent?
\item What are expectation and variance? How is the standard deviation $\sigma$ related to variance?
\item What do Bernoulli and binomial distributions describe? Explain their parameters.
\item What are joint and marginal distributions? How do you obtain a marginal distribution from a joint distribution?
\item What are covariance and correlation? Does zero covariance imply independence?
\item What is the difference between the standard deviation of an observation and the standard error of a sample mean?
\item What does the law of large numbers say about the sample mean? What additional information does the central limit theorem provide?
\item What is an estimator, and what is an estimate? Give an example distinguishing them.
\item What are bias and mean squared error? How are they related to the variance of an estimator?
\item What is a likelihood function? How does maximum likelihood estimation use it?
\item How do MLE and MAP differ? What role does the prior play in MAP estimation?
\item What is a confidence interval? What does its coverage probability mean?
\item What is the empirical distribution? How is a nonparametric bootstrap sample drawn from it?
\item What does entropy measure? What does KL divergence compare?
\item What is observed information? How is Fisher information related to it?
\end{enumerate}
}
\newpage
\blocktitle{B}{Exercises adapted from Homework 1}
These exercises separate parts of the eight homework tasks into short,
self-contained problems, with a few variations on the same models.
Throughout this handout, the second parameter of $\mathcal N(\mu,\sigma^2)$
is the variance.

\begin{question}{B1}{Inverse-CDF sampling}
A random variable has CDF
\[
F(x)=\begin{cases}
0,&x<0,\\
x^2,&0\le x\le1,\\
1,&x>1.
\end{cases}
\]
Given $U\sim\operatorname{Uniform}(0,1)$, construct a random variable with
this CDF. Find $P(X>1/2)$.
\end{question}

\begin{question}{B2}{A hidden state}
A machine is faulty with probability $0.2$. Let $Z=1$ denote a faulty
machine and $K=1$ denote a defective item. Suppose
\[
P(K=1\mid Z=0)=0.1,\qquad P(K=1\mid Z=1)=0.6.
\]
An item is defective. Find $P(Z=1\mid K=1)$.

A measurement is then recorded:
\[
Y=K+\varepsilon,\qquad \varepsilon\sim\mathcal N(0,1),
\]
where $\varepsilon$ is independent of $Z$ and $K$. Once $K$ is known,
does observing $Y$ provide additional information about $Z$? Explain.
\end{question}

\begin{question}{B3}{Waiting times and counts}
Starting at time zero, successive waiting times between requests are
independent exponential random variables with rate $2$ requests per minute.

Let $N(1)$ be the number of requests in the first minute and $T_2$ the time
of the second request. State the distributions of $N(1)$ and $T_2$.
Use the relationship between them to calculate $P(T_2\le1)$.
Leave exponentials unevaluated.
\end{question}

\begin{question}{B4}{Averaging sensor readings}
Four readings have the form
\[
X_i=B+\varepsilon_i,\qquad i=1,\ldots,4.
\]
The shared offset $B$ equals $+1$ or $-1$, each with probability $1/2$.
The noises $\varepsilon_i\sim\mathcal N(0,1)$ are independent of one another
and of $B$. Find the variance of the average reading.

Now replace the shared offset by independent offsets $B_i$, with the same
distribution and independent of all noises. What is the variance of the
average now?
\end{question}

\newpage
\blocktitle{B}{Exercises adapted from Homework 1}
\begin{question}{B5}{Correlated measurements}
Three random variables have variance $1$, and every pair has correlation
$1/2$. Calculate
\[
\Var\left(\frac{X_1+X_2+X_3}{3}\right)
\qquad\text{and}\qquad
\Var(X_1-X_2).
\]
\end{question}

\begin{question}{B6}{Distance between Gaussian vectors}
Let $X,Y$ be independent two-dimensional standard Gaussian vectors:
\[
X,Y\sim\mathcal N_2(0,I_2),
\]
where $I_2$ is the $2\times2$ identity matrix. Find the expectation and
variance of
\[
D^2=(X_1-Y_1)^2+(X_2-Y_2)^2.
\]
You may use $\E[Z^2]=1$ and $\E[Z^4]=3$ for $Z\sim\mathcal N(0,1)$.
\end{question}

\begin{question}{B7}{A Gaussian random walk}
Let $\xi_1,\ldots,\xi_5$ be independent $\mathcal N(0,1)$ variables, and define
\[
S_t=\sum_{i=1}^{t}\xi_i.
\]
Find $\Cov(S_2,S_5)$ and the distribution of $S_5-S_2$.

Is $S_5-S_2$ independent of $S_2$? Explain.
\end{question}

\begin{question}{B8}{Random linear outputs}
Let $w_1,w_2$ be independent $\mathcal N(0,1)$ variables. Two outputs use
the same weights:
\[
A=w_1+w_2,\qquad B=w_1-w_2.
\]
Find their marginal distributions and covariance. Are $A$ and $B$
independent? Justify your answer.
\end{question}

\newpage
\blocktitle{B}{Exercises adapted from Homework 1}
\begin{question}{B9}{From a CDF to moments}
A random variable takes values in $[0,1]$ and has CDF $F(x)=x^2$ on this
interval, with $F(x)=0$ below $0$ and $F(x)=1$ above $1$.

Find its density, expectation, and variance.
\end{question}

\begin{question}{B10}{An empirical CDF at one point}
Four independent observations come from a distribution on $[0,1]$ with
CDF $F(x)=x^2$ for $0\le x\le1$. The observed values are
\[
\frac14,\qquad \frac12,\qquad \frac34,\qquad \frac78.
\]
Calculate the observed fraction at or below $1/2$, and compare it with
$F(1/2)$.

If we repeatedly collect samples of size four, what are the expectation
and variance of this fraction? You may write it as
\[
\widehat F_4(1/2)=\frac14\sum_{i=1}^{4}\mathbf1\{X_i\le1/2\}.
\]
\end{question}

\begin{question}{B11}{Inference from a noisy reading}
Let $K\sim\operatorname{Bernoulli}(1/4)$ indicate whether an item is defective.
A sensor reports
\[
Y\mid K=k\sim\mathcal N(k,1).
\]
After observing $Y=1$, find $P(K=1\mid Y=1)$ and $\E[K\mid Y=1]$.
Use the Gaussian density in Bayes' formula and leave exponentials unevaluated.
\end{question}

\begin{question}{B12}{A mixture of defect rates}
A machine has one of two states, each with probability $1/2$. The number
$K$ of defective items in a batch of two satisfies
\[
K\mid Z=0\sim\operatorname{Binomial}(2,1/4),\qquad
K\mid Z=1\sim\operatorname{Binomial}(2,3/4).
\]
Find $\E[K]$ and $\Var(K)$, accounting for the randomness of the machine state.
\end{question}

\newpage
\blocktitle{B}{Exercises adapted from Homework 1}
\begin{question}{B13}{Waiting after no request has arrived}
The waiting time $T$, measured in minutes, is exponential with rate $2$.
No request has arrived during the first minute.

Find the probability of waiting at least another half-minute:
\[
P(T>1.5\mid T>1).
\]
What is the expected additional waiting time $\E[T-1\mid T>1]$?
Explain which property of the exponential distribution you use.
\end{question}

\begin{question}{B14}{Two arrival times on the same path}
Successive gaps $G_1,G_2,G_3$ are independent exponential variables with
rate $2$ requests per minute. The first and third arrival times are
\[
T_1=G_1,\qquad T_3=G_1+G_2+G_3.
\]
Find $\E[T_3]$, $\Var(T_3)$, and $\Cov(T_1,T_3)$.
State the units of the expectation and variance.
\end{question}

\begin{question}{B15}{The distribution of one sensor reading}
A sensor error is $X=B+\varepsilon$, where $B$ equals $-2$ or $2$, each
with probability $1/2$. The noise $\varepsilon\sim\mathcal N(0,1)$ is
independent of $B$.

Write the distribution of $X$ as a mixture, specifying its weights and
component distributions. Calculate $\E[X]$ and $\Var(X)$.
\end{question}

\begin{question}{B16}{Subtracting readings from the same sensor}
Two readings share an offset:
\[
X_1=B+\varepsilon_1,\qquad X_2=B+\varepsilon_2.
\]
Here $B$ equals $-1$ or $1$ with equal probability, and
$\varepsilon_1,\varepsilon_2$ are independent $\mathcal N(0,1)$ variables,
independent of $B$.

Find $\Cov(X_1,X_2)$, their correlation, and $\Var(X_1-X_2)$.
Explain what happens to the shared offset when the readings are subtracted.
\end{question}

\newpage
\blocktitle{B}{Exercises adapted from Homework 1}
\begin{question}{B17}{Can these correlations exist?}
Three zero-mean variables have variance $1$, and every distinct pair is
claimed to have correlation $-3/4$.

Calculate the implied variance of $X_1+X_2+X_3$. Can such variables exist?

Now suppose every pair has correlation $-1/2$. What does the variance of
the sum tell you about $X_1+X_2+X_3$?
\end{question}

\begin{question}{B18}{One coordinate of a random direction}
Let $X\sim\mathcal N_{100}(0,I_{100})$ and $U=X/\|X\|$, where $I_{100}$
is the $100\times100$ identity matrix. You may use that $U$ has the same
distribution after any coordinate permutation or sign change, and
\[
\sum_{i=1}^{100}U_i^2=1.
\]
Use these symmetries to find $\E[U_1]$ and $\E[U_1^2]$.
Then obtain an upper bound for $P(|U_1|\ge0.2)$ using Chebyshev's inequality.
\end{question}

\begin{question}{B19}{Overlapping random-walk increments}
Let $\xi_1,\ldots,\xi_4$ be independent $\mathcal N(0,1)$ variables and
$S_t=\sum_{i=1}^{t}\xi_i$. Consider
\[
A=S_3-S_1,\qquad B=S_4-S_2.
\]
Find $\Cov(A,B)$ and their correlation. Are $A$ and $B$ independent?
Identify the noise terms that determine your answer.
\end{question}

\begin{question}{B20}{Choosing the variance of random weights}
A fixed input is $x=(1,-1,1,-1)^\top$. A neuron has independent weights
$w_i\sim\mathcal N(0,\sigma^2)$ and output
\[
z=w^\top x.
\]
Find the distribution of $z$. Choose $\sigma^2$ so that $\Var(z)=1$.

Two neurons use independent weight vectors drawn with this variance and
the same fixed input $x$. What is the covariance of their outputs?
\end{question}

\newpage
\blocktitle{C}{New calculation problems}
\begin{question}{C1}{Reading a joint distribution}
The joint probabilities of $X,Y\in\{0,1\}$ are:
\[
\begin{array}{c|cc}
 &Y=0&Y=1\\ \hline
X=0&0.4&0.1\\
X=1&0.2&0.3
\end{array}
\]
Find $P(Y=1)$ and $P(X=1\mid Y=1)$. Are $X$ and $Y$ independent?
\end{question}

\begin{question}{C2}{A concentration bound}
Let $X_1,\ldots,X_{100}$ be independent and identically distributed, with
\[
\E[X_i]=2,\qquad \Var(X_i)=9.
\]
Use Chebyshev's inequality to obtain an upper bound for
\[
P\left(|\overline X-2|\ge0.6\right).
\]
\end{question}

\begin{question}{C3}{Approximating a sample mean}
A sample contains $64$ independent observations from the same distribution,
which has mean $5$ and variance $16$. Use the central limit theorem to
approximate
\[
P(\overline X>6).
\]
You may use $P(Z>2)=0.0228$ for a standard normal variable $Z$.
Is your answer necessarily exact under the stated assumptions?
\end{question}

\begin{question}{C4}{Method of moments}
A Gamma model has shape $\alpha$ and rate $\beta$, with
\[
\E[X]=\frac{\alpha}{\beta},\qquad
\Var(X)=\frac{\alpha}{\beta^2}.
\]
The observed sample mean is $4$, and the empirical variance
\[
\frac1n\sum_{i=1}^n(x_i-\overline x)^2
\]
equals $8$. Find the method-of-moments estimates of $\alpha$ and $\beta$.
\end{question}

\newpage
\blocktitle{C}{New calculation problems}
\begin{question}{C5}{Estimating an exponential rate}
Three independent observations from an exponential distribution with unknown
rate $\lambda$ are
\[
x_1=1,\qquad x_2=2,\qquad x_3=3.
\]
Write the likelihood and find the maximum likelihood estimate of $\lambda$.
\end{question}

\begin{question}{C6}{Adding a prior}
Six independent Bernoulli trials produce four successes. The unknown success
probability is $p$.

Find the MLE of $p$. Then find its MAP estimate under the prior density
\[
\pi(p)\propto p(1-p),\qquad 0<p<1.
\]
You do not need to calculate any normalizing constants.
\end{question}

\begin{question}{C7}{Comparing two estimators}
Two estimators of the same parameter have these properties:
\[
\begin{array}{c|cc}
\text{Estimator}&\text{Bias}&\text{Variance}\\ \hline
A&0&1/4\\
B&1/4&1/16
\end{array}
\]
Calculate their mean squared errors. Which estimator has the smaller MSE?
\end{question}

\begin{question}{C8}{A confidence interval for a mean}
A sample of $16$ independent Gaussian observations has sample mean $3$.
The population standard deviation is known to be $2$.

Construct a $95\%$ confidence interval for the population mean, using
$z_{0.975}=1.96$.

How would the interval's width change if the sample size increased to $64$,
with the same population standard deviation?
\end{question}

\newpage
\blocktitle{C}{New calculation problems}
\begin{question}{C9}{Bootstrap without a computer}
The observed dataset is $(0,2)$. A bootstrap sample consists of two
independent draws, with replacement, from these observed values.

List the possible bootstrap sample means and their probabilities.
Find the variance of the bootstrap sample mean.
\end{question}

\begin{question}{C10}{Entropy and an incorrect probability model}
The true distribution of a binary outcome is
\[
P=(1/2,1/2),
\]
while a model predicts
\[
Q=(3/4,1/4).
\]
Calculate $H(P)$, the cross-entropy $H(P,Q)$, and $D_{\mathrm{KL}}(P\|Q)$.
Use natural logarithms and leave them unevaluated.
\end{question}

\begin{question}{C11}{Curvature at the maximum}
Four independent observations follow $\mathcal N(\mu,1)$, with unknown mean
$\mu$. For the observed sample, the log-likelihood is
\[
\ell(\mu)=C-2(\mu-1)^2,
\]
where $C$ does not depend on $\mu$.

Find the MLE and the observed information
\[
J_4(\mu)=-\ell''(\mu).
\]
What is the Fisher information $I_4(\mu)$ for this model?

By how much does the log-likelihood decrease when $\mu$ moves $0.5$ away
from the MLE?
\end{question}

\begin{question}{C12}{Predicting one Gaussian variable from another}
The variables $X,Y$ are jointly Gaussian, with
\[
\E[X]=\E[Y]=0,\qquad
\Var(X)=4,\qquad
\Var(Y)=1,\qquad
\Cov(X,Y)=1.
\]
After observing $Y=2$, find the conditional expectation and conditional
variance of $X$.
\end{question}
\newpage
\blocktitle{C}{New calculation problems}
\begin{question}{C13}{Zero covariance with dependence}
The random variable $X$ takes values $-1,0,1$, each with probability $1/3$.
Let $Y=X^2$.

Calculate $\Cov(X,Y)$. Are $X$ and $Y$ independent? Support your answer
by comparing $P(Y=0\mid X=0)$ with $P(Y=0)$.
\end{question}

\begin{question}{C14}{Two estimates of a Gaussian variance}
Three independent observations from $\mathcal N(\mu,\sigma^2)$ are
\[
x_1=1,\qquad x_2=2,\qquad x_3=3.
\]
Both parameters are unknown. Calculate the MLEs of $\mu$ and $\sigma^2$.
Also calculate the usual unbiased sample variance. Explain why the two
variance estimates differ.
\end{question}

\begin{question}{C15}{A mean interval with unknown variance}
Nine independent observations from a Gaussian population have sample mean
$\overline x=4$ and sample standard deviation $s=3$, calculated using
denominator $n-1$. The population variance is unknown.

Construct an exact $95\%$ confidence interval for the population mean,
using $t_{0.975,8}=2.306$. Explain why the reference distribution is Student's
$t$ with eight degrees of freedom.
\end{question}

\begin{question}{C16}{Two bootstrap intervals}
An observed estimate is $\widehat\theta=3$. The $2.5\%$ and $97.5\%$
quantiles of its bootstrap estimates are $2$ and $5$, respectively.

Calculate the $95\%$ percentile interval and the $95\%$ basic bootstrap
interval. Why do their endpoints differ in this example?
\end{question}

\newpage
\blocktitle{C}{New calculation problems}
\begin{question}{C17}{When Markov's bound is attained}
A nonnegative random variable $X$ has expectation $3$.
Use Markov's inequality to bound
\[
P(X\ge12).
\]
Construct a distribution supported on $\{0,12\}$ with expectation $3$
for which this bound is attained.
\end{question}

\begin{question}{C18}{The score of a Bernoulli observation}
One observation $X$ has Bernoulli mass function
\[
f_p(x)=p^x(1-p)^{1-x},\qquad x\in\{0,1\}.
\]
The score is $s_p(X)=\frac{\partial}{\partial p}\log f_p(X)$.
At $p=1/4$, list its possible values and their probabilities.

Use these values to calculate $\E_p[s_p(X)]$ and the Fisher information
$I_1(p)=\E_p[s_p(X)^2]$ at $p=1/4$.
\end{question}

\begin{question}{C19}{A lower bound on estimator variance}
A sample contains $16$ independent observations from $\mathcal N(\mu,4)$.
The mean $\mu$ is unknown and the variance $4$ is known. The Fisher
information about $\mu$ in one observation is $I_1(\mu)=1/4$.

Find the information in the full sample and the Cram\'er--Rao lower bound
for the variance of an estimator that is unbiased for $\mu$ throughout
this model. Does the sample mean attain the bound? Show the calculation.
\end{question}

\begin{question}{C20}{An MLE at the boundary}
Three independent observations are $1,2,3$. The model is uniform on
$[0,\theta]$, with unknown $\theta>0$ and density
\[
f_\theta(x)=\frac1\theta\mathbf1\{0\le x\le\theta\}.
\]
Write the likelihood, including its support, and find the MLE of $\theta$.
Explain why solving $\ell'(\theta)=0$ does not find the answer.
\end{question}
\end{document}
